\section*{Capítulo \ref{ch:g7:plate-tectonics} --- Tectônica de placas: o quebra-cabeça que se move}

\begin{solution}{exo:g7:plate-tectonics:1}
A \omterm{def:g7:plate-tectonics:continental-drift}{deriva continental} é a ideia de que os \omterm{def:g3:continents-and-oceans:continent}{continentes} se deslocam devagar
pela superfície da Terra, de modo que as suas posições mudaram ao longo das
eras. As pistas de Wegener: o encaixe dos \omterm{def:g1:landscapes-around-us:water}{litorais} (a América do Sul e a
África); os mesmos fósseis em \omterm{def:g3:continents-and-oceans:continent}{continentes} hoje separados por \omterm{def:g2:water-on-earth:ocean}{oceanos}; as
mesmas \omterm{def:g3:stones-pebbles-sand:rock}{rochas} e antigas cadeias de \omterm{def:g1:landscapes-around-us:landform}{montanhas} dos dois lados de um \omterm{def:g2:water-on-earth:ocean}{oceano};
marcas de \omterm{def:g5:weather-and-climate:climate}{climas} antigos fora do lugar (carvão na Antártida, depósitos de
geleiras na África quente e árida).
\end{solution}

\begin{solution}{exo:g7:plate-tectonics:2}
Divergente: as \omterm{def:g7:plate-tectonics:plate}{placas} se afastam. Convergente: elas se aproximam.
Transformante: elas deslizam uma ao lado da outra.
\end{solution}

\begin{solution}{exo:g7:plate-tectonics:3}
Por exemplo, toda a volta do \omterm{def:g2:water-on-earth:ocean}{oceano} Pacífico (o Japão, a Indonésia, os
Andes), e a faixa que vai do Mediterrâneo até o Himalaia. Quase nenhum
grande \omterm{def:g4:volcanoes-earthquakes-first-look:earthquake}{terremoto} no interior da África ou da Austrália (nem da América do
Norte a leste das \omterm{def:g1:landscapes-around-us:landform}{Montanhas} Rochosas, nem do Brasil).
\end{solution}

\begin{solution}{exo:g7:plate-tectonics:4}
Um animal de \omterm{def:g2:water-on-earth:salt-fresh}{água doce} não conseguiria atravessar milhares de quilômetros
de \omterm{def:g2:water-on-earth:ocean}{oceano} salgado. Encontrá-lo dos dois lados significa que os dois lugares
estavam unidos por terra, num só \omterm{def:g3:continents-and-oceans:continent}{continente}, quando o \emph{Mesosaurus}
vivia.
\end{solution}

\begin{solution}{exo:g7:plate-tectonics:5}
O carvão se forma a partir das plantas de florestas quentes, úmidas e
pantanosas: o \omterm{def:g5:weather-and-climate:climate}{clima} era quente e úmido. Florestas assim não conseguem
crescer perto do \omterm{def:g3:continents-and-oceans:poles}{Polo Sul} hoje, então a Antártida deve ter ficado muito
mais perto do \omterm{def:g3:continents-and-oceans:equator}{equador}: ela se deslocou.
\end{solution}

\begin{solution}{exo:g7:plate-tectonics:6}
Dorsal mesoatlântica: divergente. Andes: convergente (uma \omterm{def:g7:plate-tectonics:plate}{placa} oceânica
afunda embaixo de um \omterm{def:g3:continents-and-oceans:continent}{continente}). Falha de San Andreas: transformante.
Himalaia: convergente (dois \omterm{def:g3:continents-and-oceans:continent}{continentes} colidem). Rifte da África
Oriental: divergente (um \omterm{def:g3:continents-and-oceans:continent}{continente} se partindo).
\end{solution}

\begin{solution}{exo:g7:plate-tectonics:7}
$2.5 \times 10 = 25$ quilômetros por milhão de anos; em 60 milhões de
anos: $25 \times 60 = \qty{1500}{km}$.
\end{solution}

\begin{solution}{exo:g7:plate-tectonics:8}
$\qty{7}{m} = \qty{700}{cm}$, e $700 \div 9 \approx 78$: cerca de
\qty{78}{cm} por ano, umas 30 vezes a média de \qty{2.5}{cm} por ano. As
\omterm{def:g7:plate-tectonics:plate}{placas} não se afastam suavemente na superfície: a tensão se acumula
durante anos e é liberada em surtos curtos, quando as fendas se abrem
metros de uma vez. Em períodos longos, os anos calmos e os surtos se
equilibram numa média de alguns centímetros por ano.
\end{solution}

\begin{solution}{exo:g7:plate-tectonics:9}
Ele não conseguia dizer que força move os \omterm{def:g3:continents-and-oceans:continent}{continentes}, e o movimento que
ele imaginava (\omterm{def:g3:continents-and-oceans:continent}{continentes} abrindo caminho através do fundo do \omterm{def:g2:water-on-earth:ocean}{oceano}) era
impossível. Na \omterm{def:g7:plate-tectonics:plate-tectonics}{tectônica de placas}, os \omterm{def:g3:continents-and-oceans:continent}{continentes} não abrem caminho
através de nada: eles são levados, com o fundo oceânico em volta deles, por
\omterm{def:g7:plate-tectonics:plate}{placas} que se movem como pedaços inteiros da casca externa da Terra.
\end{solution}

\begin{solution}{exo:g7:plate-tectonics:10}
Mais alto na dorsal, mais fundo na fossa. A \omterm{def:g7:plate-tectonics:plate}{placa} nova é criada na dorsal,
onde a \omterm{def:g3:stones-pebbles-sand:rock}{rocha} quente sobe do \omterm{def:g6:inside-the-earth:layers}{manto}; a \omterm{def:g7:plate-tectonics:plate}{placa} oceânica volta para o \omterm{def:g6:inside-the-earth:layers}{manto} na
fossa, onde ela afunda embaixo do \omterm{def:g3:continents-and-oceans:continent}{continente}.
\end{solution}

\begin{solution}{exo:g7:plate-tectonics:11}
$14.2 \div 2.2 \approx 6.45$: cerca de 6 vezes mais rápida.
Em quilômetros por milhão de anos, as velocidades são 142 e
22. Tempos para \qty{1000}{km}: $1000 \div 142 \approx 7$ milhões de anos
para a dorsal do Pacífico Leste; $1000 \div 22 \approx 45$ milhões de anos
para a dorsal do Atlântico Norte.
\end{solution}

\begin{solution}{exo:g7:plate-tectonics:12}
(a) $6472 - 4578 = 1894$ \omterm{def:g4:volcanoes-earthquakes-first-look:earthquake}{terremotos}. (b) $4578 \div 7031 \approx 0.651$,
ou seja, 65\,\%. (c) A linha no \omterm{def:g2:drawing-to-map:map}{mapa} marca a fossa, onde a \omterm{def:g7:plate-tectonics:plate}{placa} oceânica
começa a afundar. A \omterm{def:g7:plate-tectonics:plate}{placa} não desce reto: ela afunda inclinada embaixo da
outra \omterm{def:g7:plate-tectonics:plate}{placa}, cada vez mais fundo à medida que passa por baixo dela. Os
\omterm{def:g4:volcanoes-earthquakes-first-look:earthquake}{terremotos} acontecem ao longo de toda a \omterm{def:g7:plate-tectonics:plate}{placa} que afunda, então muitos
deles ficam embaixo da \omterm{def:g7:plate-tectonics:plate}{placa} de cima, do mesmo lado da fossa, a dezenas ou
centenas de quilômetros da linha. Um \omterm{def:g7:plate-tectonics:plate-boundary}{limite de placas} é uma zona, não uma
linha fina.
\end{solution}

\begin{solution}{pb:g7:plate-tectonics:1}
\textbf{1.} O \emph{Mesosaurus}, o \emph{Cynognathus} e a
\emph{Glossopteris} (a faixa do \emph{Lystrosaurus} liga a África, a Índia
e a Antártida).

\textbf{2.} Ele vivia em \omterm{def:g2:water-on-earth:salt-fresh}{água doce}: não conseguiria atravessar um \omterm{def:g2:water-on-earth:ocean}{oceano}
salgado de milhares de quilômetros de largura.

\textbf{3.} A saliência do nordeste do Brasil se encaixa no recôncavo do
golfo da Guiné, na costa oeste da África.

\textbf{4.} Um \omterm{def:g7:plate-tectonics:plate-boundary}{limite divergente}, a dorsal mesoatlântica: a \omterm{def:g7:plate-tectonics:plate}{placa}
Sul-Americana se desloca para o oeste e a \omterm{def:g7:plate-tectonics:plate}{placa} Africana para o leste, uma
se afastando da outra.

\textbf{5.} As costas ficam em lados opostos do \omterm{def:g6:latitude-longitude:longitude}{meridiano de origem}:
$43.2 + 14.4 = \ang{57.6}$.

\textbf{6.} $57.6 \times 102 = 5875.2$: cerca de \qty{5880}{km}.

\textbf{7.} Os \omterm{def:g6:latitude-longitude:longitude}{meridianos} se aproximam uns dos outros à medida que se vai
do \omterm{def:g3:continents-and-oceans:equator}{equador} para um \omterm{def:g3:continents-and-oceans:poles}{polo}. O \omterm{def:g6:latitude-longitude:latitude}{paralelo} de \ang{23} Sul é um círculo menor que
o \omterm{def:g3:continents-and-oceans:equator}{equador}, então cada um dos seus 360 graus é mais curto.

\textbf{8.} Cerca de metade da largura: $5880 \div 2 = \qty{2940}{km}$.

\textbf{9.} $3.4 \times 10 = 34$ quilômetros por milhão de anos.

\textbf{10.} $5880 \div 34 \approx 173$: cerca de 170 milhões de anos.

\textbf{11.} Sim: 170 milhões de anos fica entre 100 e 200 milhões de
anos, e vem depois do começo da fragmentação da \omterm{def:g7:plate-tectonics:supercontinent}{Pangeia}, 225 a 200 milhões
de anos atrás. Um cálculo simples num \omterm{def:g2:drawing-to-map:map}{mapa} dá a ordem de grandeza certa.

\textbf{12.} Por exemplo: a velocidade não foi sempre a mesma; os
\omterm{def:g3:continents-and-oceans:continent}{continentes} se esticaram e afinaram antes de se partir, então as costas de
hoje não são exatamente a linha ao longo da qual eles se romperam.
\end{solution}
